\section{The Supergraph Is Not a Schema}
\label{sec:ch09-not-a-schema}
\index{supergraph!what the composer writes|(}%

The command is called \code{router compose} and the thing it is understood to be
composing is a supergraph. What lands on disk is a \emph{router execution
config}: 31{,}555 bytes of JSON on a single line, holding five keys.

\begin{minted}[fontsize=\footnotesize]{text}
engineConfig            31,272 bytes
version                 "00000000-0000-0000-0000-000000000000"
subgraphs               2 entries
featureFlagConfigs      {}
compatibilityVersion    "1:0.63.2"
\end{minted}

Chapter~\ref{ch:how-a-federated-query-actually-runs} met the all-zero version on
the sample and explained it: a locally composed config has no registry to take a
version from, and the router logs those zeroes at startup as its
\code{config\_version}. Figure~\ref{fig:ch09-config-anatomy} is the whole file
in one picture, and the rest of this section is the four parts of it that carry
anything.

\begin{figure}[htbp]
  \centering
  % What one wgc router compose actually produces: three files in, one file out,
% and the four parts of that file chapter 09 takes apart. Referenced from
% chapter 09. Bare tikzpicture; the figure environment, caption and label live
% at the call site.
%
% The flow runs top to bottom because the output box is wider than any input
% and a left-to-right arrangement leaves no room for it inside the measure.
% The dashed line is routed above the input row rather than through it.
\begin{tikzpicture}[
    font=\small,
    infile/.style={draw, rounded corners=2pt, minimum width=40mm,
                   minimum height=7mm, font=\scriptsize, align=center},
    cmd/.style={draw, rounded corners=2pt, minimum width=34mm,
                minimum height=7mm, font=\scriptsize, align=center,
                fill=black!8},
    outer/.style={draw, rounded corners=2pt, minimum width=72mm,
                  minimum height=40mm},
    part/.style={draw, rounded corners=2pt, minimum width=64mm,
                 minimum height=8mm, font=\scriptsize, align=center,
                 fill=black!5},
    feed/.style={-{Stealth[length=2mm]}},
    copy/.style={-{Stealth[length=2mm]}, dashed, gray},
    lbl/.style={font=\scriptsize, align=center, inner sep=1pt}
  ]

  % -- what goes in --------------------------------------------------------
  \node[infile] (yaml) at (2.4,0)  {\code{federation/mosaic.yaml}};
  \node[infile] (cat)  at (7.0,0)  {\code{schema/catalog.graphql}};
  \node[infile] (mos)  at (11.6,0) {\code{schema/mosaic.graphql}};

  \node[cmd] (wgc) at (7.0,-1.9) {\code{wgc router compose}};

  \draw[feed] (2.4,-0.35)  -- (2.4,-1.9)  -- (wgc.west);
  \draw[feed] (7.0,-0.35)  -- (wgc.north);
  \draw[feed] (11.6,-0.35) -- (11.6,-1.9) -- (wgc.east);

  % -- what comes out ------------------------------------------------------
  \node[outer] (cfg) at (7.0,-5.15) {};
  \node[lbl, anchor=south] at (7.0,-3.05)
    {\code{federation/supergraph.json}, 31{,}555 bytes};

  \node[part] (api) at (7.0,-3.75)
    {\code{engineConfig.graphqlSchema}\\the client schema, no \code{join\_\_} in it};
  \node[part] (ds)  at (7.0,-4.75)
    {\code{datasourceConfigurations}\\which subgraph answers which fields};
  \node[part] (sdl) at (7.0,-5.75)
    {\code{federation.serviceSdl}\\each subgraph's schema, verbatim};
  \node[part] (ss)  at (7.0,-6.75)
    {\code{stringStorage}\\the same two, normalised, keyed by SHA-1};

  \draw[feed] (wgc.south) -- (7.0,-3.35);

  % The one part of the output that is a copy of the input rather than
  % something derived from it. Routed above the input row and down the right.
  \draw[copy] (cat.north) -- (7.0,0.75) -- (14.3,0.75) -- (14.3,-5.75)
              -- (sdl.east);
  \node[lbl, gray, anchor=south] at (10.65,0.8) {byte for byte};

\end{tikzpicture}

  \caption{Two schemas and a routing list in, one file out. The dashed line is
  the only part of the output that is a copy rather than a derivation: each
  subgraph's own SDL is carried through untouched. Everything else in the file
  was computed from the pair, which is why nothing in it can be edited by hand
  and still mean anything.}
  \label{fig:ch09-config-anatomy}
\end{figure}

\subsection*{The version string names two things, and both are checkable}

\index{compatibilityVersion@\code{compatibilityVersion}}%
\code{compatibilityVersion} is two numbers joined by a colon, and I chased both
because a version string nobody can decode is a version string nobody can act
on.

The \code{1} is the router contract version, and
\code{@wundergraph/composition} declares it in one line:

\begin{minted}[fontsize=\footnotesize,breaklines]{javascript}
exports.ROUTER_COMPATIBILITY_VERSION_ONE = '1';
exports.ROUTER_COMPATIBILITY_VERSIONS = new Set([exports.ROUTER_COMPATIBILITY_VERSION_ONE]);
exports.LATEST_ROUTER_COMPATIBILITY_VERSION = '1';
\end{minted}

At 0.63.2 that set has one member, so there is currently no other contract to
speak. The \code{0.63.2} is the package's own version, read from its
\code{package.json}. Between them the string says which router protocol the file
speaks and which library wrote it, and a router that refuses a config can say
which half it objected to.

That library is a dependency of \code{wgc} rather than something you install, so
its version is not one you would normally know. Everything in this chapter was
run against \code{wgc} 0.129.7, pinned exactly in the companion repository's
\code{package.json}, and this string is where the number underneath it becomes
visible.

\subsection*{Ownership is a table, not an annotation}

\index{supergraph!datasourceConfigurations@\code{datasourceConfigurations}}%
\code{engineConfig.datasourceConfigurations} is where the router learns who can
answer what. One entry per subgraph, and the interesting field on each is
\code{rootNodes}: the types the router is allowed to \emph{enter} that subgraph
at.

\begin{minted}[fontsize=\footnotesize]{text}
catalog  Query:        products, browseProducts, productById, productBySku
         Product:      id, sku, title, description, category

mosaic   Query:        customerById, ordersByCustomer
         Mutation:     submitReview
         Subscription: onReviewAdded
         Product:      availableQuantity, price, reviews, averageRating, id
\end{minted}

\index{entity!as an entry point}%
\code{Product} is a root node in both, and that is the whole of what being an
entity means, written down. Every other root node is a real operation root;
\code{Product} is a type the router can arrive at from outside, through
\code{\_entities}, carrying a key. The \code{keys} field beside it says which
key, and it is the same in both:

\begin{minted}{json}
[ { "typeName": "Product", "selectionSet": "id" } ]
\end{minted}

\code{childNodes} is the rest: five types reachable once you are inside Catalog,
fifteen inside Mosaic. Nothing surprising in either list, and the two share
\code{PageInfo}, \code{PageCursor} and \code{Node}, which is
chapter~\ref{ch:the-first-cut-extracting-catalog}'s shareable page cursor
arriving where it was always going to arrive.

\subsection*{Each subgraph's schema is in there twice}

\index{supergraph!serviceSdl@\code{serviceSdl}}%
Both copies live under \code{customGraphql}, and they are not the same text.

\medskip\noindent\textbf{\code{federation.serviceSdl} is the input, unmodified.}
I checked this rather than assuming it, because \enquote{the composer passes
your schema through} is the kind of claim that is almost true. Catalog's
committed file and the \code{serviceSdl} of datasource 0 are both 5{,}127 bytes
and hash to the same SHA-256; Mosaic's are both 5{,}150 bytes and likewise. Byte
for byte, comments and blank lines included.

\medskip\noindent\textbf{\code{upstreamSchema} is a reference to a normalised
copy.} It holds nothing but a key:

\begin{minted}{json}
"upstreamSchema": { "key": "35fef470f114894fa41bd5ab1cedf9f40b322b70" }
\end{minted}

\index{supergraph!stringStorage@\code{stringStorage}}%
That key resolves in \code{engineConfig.stringStorage}, which here has exactly
two entries, one per subgraph. And the key of each entry is the SHA-1 of its own
value. I hashed both to be sure, and both match. So the config is content
addressed, and a graph with a hundred subgraphs in it stores one copy per
distinct schema rather than one per subgraph. Mosaic has two subgraphs and no
two identical schemas, so I have not seen that case; it is what content
addressing is for rather than something this file demonstrates.

The normalised form is 4{,}224 bytes against the 5{,}127 it came from, and the
900 bytes it lost are the federation machinery. \code{Query} in the normalised
copy has no \code{\_service} and no \code{\_entities} on it, and \code{\_Service},
\code{\_Entity} and \code{\_Any} are not in the document at all. In their place
are the definitions of \code{@key}, \code{@link} and \code{@shareable}, plus
\code{link\_\_Import}, \code{link\_\_Purpose} and \code{openfed\_\_FieldSet}.

\index{supergraph!normalisation}%
Which is worth pausing on, because it says what the copy is for. The router does
not need to be told that a subgraph has an \code{\_entities} field; every
subgraph has one, and the router is going to call it. What it needs is the
declarations that let it check a query against that subgraph, and the
normalisation keeps exactly those.

The rest of the rewriting is presentational, and the point of it is that two
schemas which mean the same thing end up spelled the same way. Types, fields and
arguments are all sorted alphabetically. Descriptions become block strings. The
multi-line \code{@link} collapses onto one line. \code{@key(fields: FieldSet!)}
becomes \code{openfed\_\_FieldSet!}. And an object-level \code{@shareable} is
pushed down onto every field of the object, so
chapter~\ref{ch:the-first-cut-extracting-catalog}'s

\begin{graphqlsdl}
type PageCursor @shareable {
  page: Int!
  cursor: String!
}
\end{graphqlsdl}

arrives as

\begin{graphqlsdl}
type PageCursor {
  cursor: String! @shareable
  page: Int! @shareable
}
\end{graphqlsdl}

\index{listings!descriptions elided}%
Two things about how the SDL listings in this chapter are printed, since one of
them is on the page above. Description strings are taken out, because the real
text carries a docstring above almost every field and printing them would triple
every listing. And where a field's arguments are one per line only because each
one has a description, the arguments are closed back up onto one line. Nothing
else is altered, and every file is in the repository if you want it untouched.

Same meaning either way, one canonical spelling of it. A composer that has to
compare two schemas has to decide what \enquote{the same} means before it can
compare anything, and this is that decision made concrete.

\subsection*{No join directives anywhere}

\index{supergraph!join directives@\code{join\_\_} directives}%
\code{engineConfig.graphqlSchema} is the merged schema, 6{,}035 characters of
SDL as one string. It is also where I expected to find the routing information
and did not.

Apollo describes three schemas. A subgraph schema, a supergraph schema that
\enquote{combines all of the types and fields from your subgraph schemas, plus
some federation-specific information that tells your router which subgraphs can
resolve which fields}, and an API schema that \enquote{omits federation-specific
types, fields, and directives that are considered
\enquote{machinery}}~\autocite{apollo2026schematypes}. In Apollo's stack the
middle one is a real document you can read, annotated throughout with
\code{@join\_\_type}, \code{@join\_\_field} and \code{@join\_\_graph}.

Grep the entire 31 KB Cosmo config for \code{join\_\_} and you get nothing.
There is no annotated middle document in it at all. What Apollo puts in one
document, Cosmo splits in two: a clean client-facing schema, and the
\code{datasourceConfigurations} table above.

Here is \code{Product} as the merged schema has it:

\begin{graphqlsdl}
type Product implements Node {
  id: ID!
  sku: String!
  title: String!
  description: String
  category: ProductCategory!
  availableQuantity: Int!
  price: Money!
  reviews(first: Int, after: String, last: Int, before: String): ReviewConnection!
  averageRating: Float
}
\end{graphqlsdl}

Nine fields from two services, and nothing on the page says so. The only
directive in the whole merged schema is the \code{@deprecated} that
chapter~\ref{ch:schema-design-that-survives-change} put on \code{products}.

\index{Node interface@\code{Node} interface!after composition}%
\code{implements Node} survived, and \code{Query.node} did not, which is
chapter~\ref{ch:the-first-cut-extracting-catalog}'s decision visible from the
outside for the first time: the interface is still in the graph, and the root
field you would hand an identifier back to is gone.

I do not think either arrangement is better and I would not choose a stack on
it. What it changes is where you look. If you are used to reading an Apollo
supergraph to find out which service owns a field, the equivalent question here
is answered by a JSON array rather than by a directive, and no amount of reading
the schema will answer it.

\subsection*{One type in there that nobody wrote}

\index{supergraph!FieldSet leak@\code{FieldSet} leak}%
The merged schema contains this:

\begin{graphqlsdl}
scalar FieldSet
\end{graphqlsdl}

Nothing in that schema uses it. No field returns it, no argument takes it, and a
client has no way to send one.

It arrives because of how HotChocolate writes its \code{@link}:

\begin{graphqlsdl}
schema
  @link(
    url: "https://specs.apollo.dev/federation/v2.6"
    import: ["@key", "@shareable", "@tag", "FieldSet"]
  )
\end{graphqlsdl}

\code{FieldSet} is imported as a named type rather than only as the type of
\code{@key}'s argument, so the composer treats it as a type the subgraph
declared and keeps it. Both subgraphs do it, so it survives the merge.

It is harmless and it is visible to anyone who introspects the router, which
makes it the second thing in two chapters where HotChocolate's output and the
Cosmo composer disagree about what belongs in a federated schema. The first was
the \code{@cost} directive of
chapter~\ref{ch:the-first-cut-extracting-catalog}, and I declined to say which
tool was wrong then for the same reason I decline now: two implementations
reading one specification differently is a fact about your build, and picking a
side does not change what ships.

Everything so far is what the composer produces when it is content. The next
section is the one check that has no equivalent in a single service, and the
only way to see it is to break the graph on purpose.

\index{supergraph!what the composer writes|)}%
